\documentclass[11pt]{article}

\usepackage[a4paper,margin=1in]{geometry}
\usepackage{amsmath,amssymb,amsthm,mathtools}
\usepackage{bm}
\usepackage{booktabs}
\usepackage{array}
\usepackage{enumitem}
\usepackage{hyperref}
\usepackage{xcolor}

\usepackage{booktabs}
\usepackage{multirow}
\usepackage{array}
\usepackage{makecell}

\title{Understanding Search vs.\ Decision vs.\ Distinguishing\\
	in Post-Quantum Cryptography Hard Problems}
\author{}
\date{}

\newtheorem{definition}{Definition}
\newtheorem{remark}{Remark}
\newtheorem{example}{Example}

\begin{document}
	\maketitle
	
	\section{Why this distinction matters}
	
	In post-quantum cryptography (PQC), many security assumptions are stated using one of three flavors of hard problems:
	
	\begin{enumerate}[label=\arabic*.]
		\item \textbf{Search problem:} recover the hidden object.
		\item \textbf{Decision problem:} answer a yes/no question about the instance.
		\item \textbf{Distinguishing problem:} tell whether samples come from one distribution or another.
	\end{enumerate}
	
	These three look similar, but they play different roles in theory and in cryptographic reductions.
	
	\medskip
	
	\noindent
	Roughly:
	\[
	\text{Search} \quad \Rightarrow \quad \text{Decision} \quad \Rightarrow \quad \text{Distinguishing}
	\]
	in the sense that if you can fully recover the secret, then you can usually answer the weaker yes/no or distinguishing task. The converse is usually harder and may require nontrivial reductions.
	
	\section{The core intuition}
	
	Suppose Nature hides some secret structure.
	
	\begin{itemize}
		\item In a \textbf{search} problem, you must \emph{find} that hidden structure.
		\item In a \textbf{decision} problem, you only need to \emph{decide whether a certain property holds}.
		\item In a \textbf{distinguishing} problem, you do not need the secret or even a property of one fixed instance; you only need to \emph{tell apart two candidate generating processes}.
	\end{itemize}
	
	So the required output becomes progressively weaker.
	
	\section{Formal viewpoint}
	
	\subsection{Search problems}
	
	\begin{definition}[Search problem]
		A search problem consists of an instance space $\mathcal{I}$ and a relation
		\[
		R \subseteq \mathcal{I} \times \mathcal{W},
		\]
		where $\mathcal{W}$ is the witness or solution space. Given an instance $I \in \mathcal{I}$, the goal is to output some $w \in \mathcal{W}$ such that
		\[
		(I,w)\in R.
		\]
	\end{definition}
	
	\noindent
	The output is not just yes/no. It must be an actual solution.
	
	\begin{example}[Search-LWE]
		Given samples
		\[
		(\bm a_i,\, b_i = \langle \bm a_i,\bm s\rangle + e_i \bmod q),
		\]
		recover the secret vector $\bm s$.
	\end{example}
	
	\begin{example}[SIS as search]
		Given a matrix $A \in \mathbb{Z}_q^{m\times n}$, find a nonzero short vector $\bm z \in \mathbb{Z}^n$ such that
		\[
		A\bm z \equiv \bm 0 \pmod q.
		\]
	\end{example}
	
	\subsection{Decision problems}
	
	\begin{definition}[Decision problem]
		A decision problem is specified by a language $L \subseteq \mathcal{I}$. Given an instance $I$, the goal is to output
		\[
		\begin{cases}
			1 & \text{if } I \in L,\\
			0 & \text{if } I \notin L.
		\end{cases}
		\]
	\end{definition}
	
	\noindent
	The output is only one bit.
	
	\begin{example}[GapSVP]
		Given a lattice basis $B$ and threshold $d$, decide whether the shortest nonzero lattice vector is
		\[
		\lambda_1(\mathcal{L}(B)) \le d
		\]
		or
		\[
		\lambda_1(\mathcal{L}(B)) > \gamma d,
		\]
		for some approximation factor $\gamma > 1$, assuming the input is promised to satisfy one of these two cases.
	\end{example}
	
	\begin{remark}
		This is not a search problem because you do \emph{not} need to output the short vector itself. You only need to decide which case holds.
	\end{remark}
	
	\subsection{Distinguishing problems}
	
	\begin{definition}[Distinguishing problem]
		A distinguishing problem gives oracle access or samples from one of two distributions
		\[
		\mathcal{D}_0,\mathcal{D}_1.
		\]
		The goal is to decide from which distribution the samples were drawn, with non-negligible advantage over random guessing.
	\end{definition}
	
	\noindent
	This is often written as:
	\[
	\left| \Pr[\mathcal{A}(\mathcal{D}_0)=1] - \Pr[\mathcal{A}(\mathcal{D}_1)=1] \right|
	\]
	being non-negligible.
	
	\begin{example}[Decision-LWE as a distinguishing problem]
		Given samples, distinguish between
		\[
		(\bm a,\, \langle \bm a,\bm s\rangle + e \bmod q)
		\]
		and
		\[
		(\bm a,\, u),
		\]
		where $u \leftarrow \mathbb{Z}_q$ is uniform.
	\end{example}
	
	\begin{remark}
		In modern cryptography, many so-called \textbf{decision assumptions} are really \textbf{distribution distinguishing assumptions}. Decision-LWE is the canonical example.
	\end{remark}
	
	\section{A clean conceptual separation}
	
	\subsection{Search asks ``What is the hidden object?''}
	
	The hidden object may be:
	\begin{itemize}
		\item a secret vector $\bm s$ in LWE,
		\item a short relation $\bm z$ in SIS,
		\item a short lattice vector in SVP,
		\item a close lattice vector in CVP.
	\end{itemize}
	
	\noindent
	The algorithm must actually produce it.
	
	\subsection{Decision asks ``Does this instance have the property?''}
	
	Examples:
	\begin{itemize}
		\item Is there a vector shorter than threshold $d$?
		\item Is this lattice instance in the YES-case or NO-case?
		\item Is this polynomial system solvable under some promise?
	\end{itemize}
	
	\subsection{Distinguishing asks ``Which world am I in?''}
	
	You are shown samples from one of two worlds:
	\begin{itemize}
		\item \textbf{structured world} (contains hidden algebraic or noisy linear structure),
		\item \textbf{random world} (uniform or pseudorandom noise only).
	\end{itemize}
	
	\noindent
	The task is to identify the world, not necessarily to recover the structure.
	
	\section{The standard cryptographic examples}
	
	\subsection{LWE}
	
	Let $\bm s \in \mathbb{Z}_q^n$ be secret. Samples are of the form
	\[
	(\bm a, b), \qquad b = \langle \bm a,\bm s\rangle + e \pmod q,
	\]
	where $\bm a \leftarrow \mathbb{Z}_q^n$ uniformly and $e$ is drawn from some error distribution.
	
	\paragraph{Search-LWE.}
	Given many samples $(\bm a_i,b_i)$, recover $\bm s$.
	
	\paragraph{Decision-LWE / Distinguishing-LWE.}
	Given many samples, decide whether they are:
	\begin{itemize}
		\item valid LWE samples with a fixed secret $\bm s$, or
		\item uniformly random samples in $\mathbb{Z}_q^n \times \mathbb{Z}_q$.
	\end{itemize}
	
	\paragraph{Interpretation.}
	\begin{itemize}
		\item Search-LWE asks for the \emph{secret}.
		\item Decision-LWE asks whether the sample set is \emph{structured} or \emph{random}.
	\end{itemize}
	
	\subsection{SIS}
	
	Let $A \in \mathbb{Z}_q^{m\times n}$.
	
	\paragraph{Search-SIS.}
	Find a short nonzero vector $\bm z \in \mathbb{Z}^n$ such that
	\[
	A\bm z \equiv \bm 0 \pmod q.
	\]
	
	\paragraph{Decision analogue?}
	SIS is usually used primarily as a \emph{search} assumption. A natural decision-style version would be something like:
	\begin{quote}
		Does there exist a nonzero vector $\bm z$ with $\|\bm z\| \le \beta$ such that $A\bm z \equiv \bm 0 \pmod q$?
	\end{quote}
	But in cryptography, SIS is most often stated and used in its search form.
	
	\subsection{Lattice geometry problems}
	
	\paragraph{SVP (Shortest Vector Problem).}
	Given a lattice basis $B$, output a shortest nonzero vector in $\mathcal{L}(B)$.
	
	This is a \textbf{search} problem.
	
	\paragraph{GapSVP.}
	Given $B$ and threshold $d$, decide whether
	\[
	\lambda_1(\mathcal{L}(B)) \le d
	\]
	or
	\[
	\lambda_1(\mathcal{L}(B)) > \gamma d.
	\]
	
	This is a \textbf{decision} problem with a promise gap.
	
	\paragraph{CVP (Closest Vector Problem).}
	Given basis $B$ and target $\bm t$, output a lattice vector closest to $\bm t$.
	
	This is a \textbf{search} problem.
	
	\paragraph{BDD (Bounded Distance Decoding).}
	Given a target known to be unusually close to the lattice, recover the nearby lattice point.
	
	This is also fundamentally a \textbf{search} problem.
	
	\section{Why distinguishing is not exactly the same as ordinary decision}
	
	Students often say:
	
	\begin{quote}
		``A distinguishing problem is just a decision problem, because the answer is yes/no.''
	\end{quote}
	
	This is \emph{partly true} but conceptually incomplete.
	
	\subsection{Common ground}
	
	Both decision and distinguishing output one bit. So at the most abstract complexity-theoretic level, a distinguisher computes a Boolean function.
	
	\subsection{Important difference}
	
	In a \textbf{decision problem}, the input is usually a \emph{single instance} and you ask whether it belongs to a language.
	
	In a \textbf{distinguishing problem}, the input is typically generated from one of \emph{two probability distributions}. The challenge is statistical or computational separation of distributions.
	
	\medskip
	
	\noindent
	So:
	
	\begin{itemize}
		\item \textbf{Decision} is language-based.
		\item \textbf{Distinguishing} is distribution-based.
	\end{itemize}
	
	\noindent
	In cryptography, the distribution-based viewpoint is often the more natural one.
	
	\section{Reduction intuition between the three}
	
	\subsection{From search to distinguishing}
	
	If you can solve the search problem, you can often distinguish.
	
	\begin{example}[Search-LWE $\Rightarrow$ Distinguishing-LWE]
		Suppose you can recover $\bm s$. Then, on input samples $(\bm a_i,b_i)$:
		\begin{enumerate}
			\item Recover candidate $\bm s$.
			\item Check whether each
			\[
			b_i - \langle \bm a_i,\bm s\rangle \pmod q
			\]
			looks like small noise.
		\end{enumerate}
		If yes, declare ``LWE''; otherwise declare ``uniform''.
	\end{example}
	
	\noindent
	So search is at least as strong as distinguishing in this sense.
	
	\subsection{From distinguishing to search}
	
	This direction is much subtler. One usually needs a theorem stating that a distinguisher can be transformed into a search solver.
	
	For LWE, there are standard reductions showing that under appropriate parameters, \emph{decision-LWE and search-LWE are polynomial-time equivalent}.
	
	\begin{remark}
		This equivalence is nontrivial. It does not follow merely from intuition; it requires a formal reduction.
	\end{remark}
	
	\subsection{From search to ordinary decision}
	
	If you can find the exact witness, then you can certainly answer whether one exists or whether the relevant threshold is met.
	
	For example:
	\[
	\text{SVP solver} \Rightarrow \text{GapSVP solver}.
	\]
	
	But the reverse need not hold in any obvious way, because knowing only that a short vector exists does not directly tell you what the vector is.
	
	\section{A geometric picture}
	
	\subsection{Search in lattices}
	
	You are standing outside a lattice and must point to a specific geometric object:
	\begin{itemize}
		\item the shortest nonzero lattice vector,
		\item the lattice point closest to a target,
		\item a short modular dependency.
	\end{itemize}
	
	\subsection{Decision in lattices}
	
	You do not need the vector. You only answer:
	\begin{quote}
		``Is there a vector this short?''
	\end{quote}
	
	\subsection{Distinguishing in noisy linear algebra}
	
	You are shown clouds of samples:
	\begin{itemize}
		\item one cloud comes from a hidden affine/noisy structure,
		\item the other cloud is uniform noise.
	\end{itemize}
	You only decide which cloud-generating mechanism is more likely.
	
	\section{Practical role in cryptography}
	
	\subsection{Search assumptions}
	
	Search assumptions are common when the cryptographic break itself requires producing a concrete object.
	
	Examples:
	\begin{itemize}
		\item forging a signature may amount to finding a short SIS solution,
		\item decoding may amount to finding a close vector,
		\item key recovery is a search task.
	\end{itemize}
	
	\subsection{Distinguishing assumptions}
	
	Distinguishing assumptions are common in encryption and pseudorandomness.
	
	Examples:
	\begin{itemize}
		\item IND-CPA security is itself a distinguishing game,
		\item LWE naturally gives pseudorandom-looking samples,
		\item module-LWE and ring-LWE assumptions are usually phrased as distinguishing assumptions.
	\end{itemize}
	
	\subsection{Decision assumptions}
	
	Decision versions often appear in complexity theory and worst-case to average-case reductions.
	
	Examples:
	\begin{itemize}
		\item GapSVP,
		\item GapCVP,
		\item approximate threshold formulations used in hardness proofs.
	\end{itemize}
	
	\section{Summary table}
	
	\begin{center}
		\renewcommand{\arraystretch}{1.35}
		\begin{tabular}{>{\raggedright\arraybackslash}p{2.5cm} >{\raggedright\arraybackslash}p{4.1cm} >{\raggedright\arraybackslash}p{4.5cm} >{\raggedright\arraybackslash}p{3cm}}
			\toprule
			\textbf{Type} & \textbf{Output} & \textbf{Typical question} & \textbf{PQC examples} \\
			\midrule
			Search & explicit witness / secret / vector & ``Find the hidden object'' & Search-LWE, SIS, SVP, CVP, BDD \\
			\midrule
			Decision & one bit about instance property & ``Does this property hold?'' & GapSVP, GapCVP, threshold versions of lattice problems \\
			\midrule
			Distinguishing & one bit about source distribution & ``Are these samples structured or random?'' & Decision-LWE, Ring-LWE, Module-LWE \\
			\bottomrule
		\end{tabular}
	\end{center}
	
	\section{A very important warning}
	
	In PQC papers, the word \textbf{decision} is sometimes used loosely.
	
	\begin{itemize}
		\item In classical complexity theory, ``decision'' usually means language membership.
		\item In cryptography, ``decision-LWE'' usually means distinguishing LWE samples from uniform.
	\end{itemize}
	
	So when reading a paper, always ask:
	
	\begin{quote}
		Is ``decision'' here an \emph{instance property problem}, or a \emph{distribution distinguishing problem}?
	\end{quote}
	
	\section{One-line memory aid}
	
	\[
	\boxed{
		\begin{aligned}
			\text{Search} &:\ \text{recover the object},\\
			\text{Decision} &:\ \text{decide the property},\\
			\text{Distinguishing} &:\ \text{identify the distribution}.
		\end{aligned}
	}
	\]
	
	\section{Mini-exercises}
	
	\begin{enumerate}
		\item Explain why recovering the LWE secret immediately gives a distinguisher.
		\item Explain why GapSVP is a decision problem but SVP is a search problem.
		\item For SIS, formulate a natural yes/no version and compare it with the usual search version.
		\item In one sentence each, state the difference between:
		\begin{itemize}
			\item Search-LWE and Decision-LWE,
			\item SVP and GapSVP,
			\item decision and distinguishing.
		\end{itemize}
	\end{enumerate}
	
	\section{Final takeaway}
	
	The cleanest way to think about PQC hard problems is:
	
	\begin{itemize}
		\item \textbf{Search:} produce the hidden mathematical object.
		\item \textbf{Decision:} answer a yes/no structural question about one instance.
		\item \textbf{Distinguishing:} tell whether observed data comes from a structured source or a random source.
	\end{itemize}
	
	\noindent
	For lattice-based cryptography in particular:
	\begin{itemize}
		\item \textbf{SVP/CVP/SIS} are naturally search-oriented,
		\item \textbf{GapSVP/GapCVP} are naturally decision-oriented,
		\item \textbf{LWE/RLWE/MLWE} are most often used in distinguishing form.
	\end{itemize}

	\newpage
	\section{Code-based hard problems: search vs.\ decision vs.\ distinguishing}
	
	Code-based cryptography gives another clean setting in which the distinction between search, decision, and distinguishing becomes very important.
	
	\subsection{Basic coding-theory setup}
	
	Let $\mathbb{F}_q$ be a finite field, often $\mathbb{F}_2$ in classical code-based cryptography.
	
	A linear code $C$ of length $n$ and dimension $k$ is a $k$-dimensional subspace
	\[
	C \subseteq \mathbb{F}_q^n.
	\]
	
	It is commonly represented by either:
	
	\begin{itemize}
		\item a \textbf{generator matrix} $G \in \mathbb{F}_q^{k \times n}$ such that
		\[
		C = \{ \bm m G : \bm m \in \mathbb{F}_q^k \},
		\]
		\item or a \textbf{parity-check matrix} $H \in \mathbb{F}_q^{(n-k)\times n}$ such that
		\[
		C = \{ \bm c \in \mathbb{F}_q^n : H\bm c^T = \bm 0 \}.
		\]
	\end{itemize}
	
	The \emph{Hamming weight} of a vector $\bm e \in \mathbb{F}_q^n$, written $w_H(\bm e)$, is the number of nonzero coordinates.
	
	\medskip
	
	\noindent
	In code-based cryptography, the hidden object is often a \textbf{low-weight error vector}.
	
	\subsection{The central problem: syndrome decoding}
	
	Given a parity-check matrix $H$ and a syndrome $\bm s$, the fundamental equation is
	\[
	H \bm e^T = \bm s^T.
	\]
	
	The unknown is the error vector $\bm e$.
	
	\paragraph{Why this matters.}
	If $\bm y = \bm c + \bm e$ is a noisy codeword, then
	\[
	H \bm y^T = H(\bm c+\bm e)^T = H\bm c^T + H\bm e^T = \bm 0 + H\bm e^T = \bm s^T.
	\]
	So the syndrome depends only on the error pattern, not on the codeword itself.
	
	\subsection{Search version: Syndrome Decoding (SD)}
	
	\begin{definition}[Search Syndrome Decoding]
		Given
		\[
		H \in \mathbb{F}_q^{(n-k)\times n}, \qquad \bm s \in \mathbb{F}_q^{n-k}, \qquad w,
		\]
		find a vector $\bm e \in \mathbb{F}_q^n$ such that
		\[
		H\bm e^T = \bm s^T
		\qquad \text{and} \qquad
		w_H(\bm e) \le w.
		\]
	\end{definition}
	
	This is a \textbf{search problem}: you must output the actual low-weight error vector.
	
	\begin{remark}
		This is the code-based analogue of many PQC search problems:
		\begin{itemize}
			\item in LWE search, recover the secret;
			\item in SIS, find a short relation;
			\item in syndrome decoding, find the sparse error.
		\end{itemize}
	\end{remark}
	
	\subsection{Decision version}
	
	A natural decision analogue is:
	
	\begin{definition}[Decision Syndrome Decoding]
		Given
		\[
		H,\bm s,w,
		\]
		decide whether there exists some vector $\bm e$ such that
		\[
		H\bm e^T = \bm s^T
		\qquad \text{and} \qquad
		w_H(\bm e) \le w.
		\]
	\end{definition}
	
	Here the output is only:
	\[
	\text{YES or NO.}
	\]
	
	\noindent
	You do not need to produce the error vector.
	
	\subsection{Gap-style decision version}
	
	Just as lattices have GapSVP and GapCVP, code problems can also be phrased in threshold or promise-gap form.
	
	\begin{definition}[Gap decoding style problem]
		Given $H,\bm s,w$ and approximation factor $\gamma > 1$, distinguish between:
		\begin{itemize}
			\item \textbf{YES-case:} there exists $\bm e$ with
			\[
			H\bm e^T = \bm s^T
			\quad \text{and} \quad
			w_H(\bm e) \le w,
			\]
			\item \textbf{NO-case:} every $\bm e$ satisfying
			\[
			H\bm e^T = \bm s^T
			\]
			has
			\[
			w_H(\bm e) > \gamma w.
			\]
		\end{itemize}
	\end{definition}
	
	This is conceptually parallel to lattice gap problems.
	
	\subsection{Distinguishing version}
	
	In code-based cryptography, distinguishing problems often ask whether some public object comes from a \emph{structured code family} or from a \emph{random linear code}.
	
	\paragraph{Generic form.}
	Given some matrix or samples, distinguish between:
	
	\begin{itemize}
		\item \textbf{structured world:} generated from a code with hidden algebraic structure,
		\item \textbf{random world:} generated uniformly or from an unstructured random code ensemble.
	\end{itemize}
	
	\subsection{Example: McEliece-style viewpoint}
	
	In the McEliece framework, one starts with a secret code that has efficient decoding structure, then disguises it publicly.
	
	Very roughly:
	\begin{itemize}
		\item the \textbf{secret key} contains a highly structured code with an efficient decoder,
		\item the \textbf{public key} looks like a generator matrix of a random linear code.
	\end{itemize}
	
	A natural distinguishing question is:
	
	\begin{quote}
		Can an adversary distinguish the public code from a uniformly random linear code?
	\end{quote}
	
	That is a \textbf{distinguishing problem}, not a search problem.
	
	\begin{remark}
		This is different from syndrome decoding:
		\begin{itemize}
			\item \textbf{Syndrome decoding} asks to recover a low-weight error;
			\item \textbf{code indistinguishability} asks whether the code itself hides special algebraic structure.
		\end{itemize}
	\end{remark}
	
	\subsection{Geometric intuition in Hamming space}
	
	Lattice problems live in Euclidean geometry.  
	Code-based problems live in \textbf{Hamming geometry}.
	
	\begin{itemize}
		\item A sparse error vector means a vector with few nonzero coordinates.
		\item Decoding asks for a codeword close to the received word in Hamming distance.
		\item Syndrome decoding asks for the sparse error pattern that explains the syndrome.
	\end{itemize}
	
	So the analogies are:
	
	\[
	\begin{aligned}
		\text{short vector in lattices}
		&\longleftrightarrow
		\text{low-weight vector in codes},\\[4pt]
		\text{Euclidean distance}
		&\longleftrightarrow
		\text{Hamming distance},\\[4pt]
		\text{closest vector problem}
		&\longleftrightarrow
		\text{nearest codeword / decoding problem}.
	\end{aligned}
	\]
	
	\subsection{Search vs.\ decision vs.\ distinguishing in code-based PQC}
	
	\paragraph{Search.}
	Recover the hidden sparse object.
	
	Typical question:
	\begin{quote}
		Find $\bm e$ such that $H\bm e^T = \bm s^T$ and $w_H(\bm e)\le w$.
	\end{quote}
	
	\paragraph{Decision.}
	Determine whether such a sparse object exists.
	
	Typical question:
	\begin{quote}
		Does there exist an error vector of weight at most $w$ consistent with this syndrome?
	\end{quote}
	
	\paragraph{Distinguishing.}
	Tell structured-code distributions apart from random-code distributions.
	
	Typical question:
	\begin{quote}
		Does this public matrix come from a hidden Goppa-like construction, or is it just a random linear code?
	\end{quote}
	
	\subsection{Why the distinction matters in cryptography}
	
	These three forms appear in different security roles.
	
	\begin{itemize}
		\item \textbf{Search hardness} is central when breaking the scheme means recovering an error pattern or decoding.
		\item \textbf{Decision hardness} is useful for complexity-theoretic classification and threshold formulations.
		\item \textbf{Distinguishing hardness} is relevant when public keys are supposed to look random and leak no visible structure.
	\end{itemize}
	
	\subsection{Comparison with lattice-based assumptions}
	
	The analogy with lattice-based cryptography is very useful:
	
	\begin{center}
		\renewcommand{\arraystretch}{1.35}
		\begin{tabular}{>{\raggedright\arraybackslash}p{3.3cm} >{\raggedright\arraybackslash}p{5cm} >{\raggedright\arraybackslash}p{5cm}}
			\toprule
			\textbf{Category} & \textbf{Lattice-based} & \textbf{Code-based} \\
			\midrule
			Search & SVP, CVP, SIS, Search-LWE & Syndrome Decoding, low-weight codeword search, decoding \\
			\midrule
			Decision & GapSVP, GapCVP & existence of error of weight $\le w$, gap decoding formulations \\
			\midrule
			Distinguishing & Decision-LWE, Ring-LWE, Module-LWE & structured code vs.\ random code, public-key indistinguishability \\
			\bottomrule
		\end{tabular}
	\end{center}
	
	\subsection{A very important caution}
	
	In code-based cryptography, not every assumption is naturally a distinguishing assumption.
	
	For example:
	\begin{itemize}
		\item \textbf{syndrome decoding} is fundamentally a \emph{search} problem;
		\item \textbf{McEliece key indistinguishability} is fundamentally a \emph{distinguishing} problem.
	\end{itemize}
	
	So one should not compress all code-based hardness into a single category.
	
	\subsection{One-line memory aid for code-based PQC}
	
	\[
	\boxed{
		\begin{aligned}
			\text{Search} &:\ \text{find the low-weight error},\\
			\text{Decision} &:\ \text{decide whether such an error exists},\\
			\text{Distinguishing} &:\ \text{tell structured codes from random codes.}
		\end{aligned}
	}
	\]
	
	\subsection{Mini-exercises on code-based problems}
	
	\begin{enumerate}
		\item Explain why recovering the error vector in syndrome decoding immediately solves the corresponding decision problem.
		\item State the difference between:
		\begin{itemize}
			\item decoding a received word,
			\item syndrome decoding,
			\item distinguishing a structured code from a random code.
		\end{itemize}
		\item Give the Hamming-space analogue of:
		\begin{itemize}
			\item a short lattice vector,
			\item a closest lattice vector,
			\item a gap decision problem.
		\end{itemize}
	\end{enumerate}

	\newpage
	\section{Summary table across PQC families}
	
	\begin{center}
		\renewcommand{\arraystretch}{1.35}
		\begin{tabular}{>{\raggedright\arraybackslash}p{2.3cm} >{\raggedright\arraybackslash}p{3.8cm} >{\raggedright\arraybackslash}p{4.6cm} >{\raggedright\arraybackslash}p{4.2cm}}
			\toprule
			\textbf{Type} & \textbf{Output} & \textbf{Lattice-based examples} & \textbf{Code-based examples} \\
			\midrule
			Search & explicit witness / secret / sparse object & SVP, CVP, BDD, SIS, Search-LWE & Syndrome Decoding, decoding, low-weight codeword search \\
			\midrule
			Decision & one bit about instance property & GapSVP, GapCVP & existence of error of weight $\le w$, gap decoding forms \\
			\midrule
			Distinguishing & one bit about source distribution & Decision-LWE, Ring-LWE, Module-LWE & structured code vs.\ random code, disguised secret code vs.\ random linear code \\
			\bottomrule
		\end{tabular}
	\end{center}

	\newpage
	\section{Final takeaway}
	
	The cleanest way to think about PQC hard problems is:
	
	\begin{itemize}
		\item \textbf{Search:} produce the hidden mathematical object.
		\item \textbf{Decision:} answer a yes/no structural question about one instance.
		\item \textbf{Distinguishing:} tell whether observed data comes from a structured source or a random source.
	\end{itemize}
	
	\noindent
	Across major PQC families:
	\begin{itemize}
		\item \textbf{lattice-based:} SVP/CVP/SIS are search-oriented, GapSVP is decision-oriented, and LWE-type assumptions are often distinguishing-oriented;
		\item \textbf{code-based:} syndrome decoding is search-oriented, threshold decoding variants are decision-oriented, and code indistinguishability questions are distinguishing-oriented.
	\end{itemize}

	\newpage
	\section{Compact summary table}
	
	\section{PQC hard problems classified as Search, Decision, and Distinguishing}
	
	\begin{center}
		\scriptsize
		\renewcommand{\arraystretch}{1.35}
		\begin{tabular}{|p{1.9cm}|p{2.9cm}|p{3.3cm}|p{4.4cm}|p{4.8cm}|}
			\hline
			\textbf{PQC family} & \textbf{Representative hard problem / assumption} & \textbf{Type} & \textbf{Informal task} & \textbf{Typical formulation} \\
			\hline
			
			\textbf{Lattice-based} 
			& SVP (Shortest Vector Problem) 
			& Search 
			& Find the shortest nonzero lattice vector 
			& Given a basis $B$ of a lattice $\mathcal{L}(B)$, output $\bm v \in \mathcal{L}(B)\setminus\{\bm 0\}$ such that $\|\bm v\|=\lambda_1(\mathcal{L}(B))$. \\
			\hline
			
			\textbf{Lattice-based} 
			& Approx-SVP 
			& Search 
			& Find a sufficiently short nonzero lattice vector 
			& Given $B$, output $\bm v \neq \bm 0$ with $\|\bm v\| \le \gamma \lambda_1(\mathcal{L}(B))$ for approximation factor $\gamma \ge 1$. \\
			\hline
			
			\textbf{Lattice-based} 
			& GapSVP 
			& Decision 
			& Decide whether the lattice has a very short vector 
			& Given $(B,d)$, distinguish between $\lambda_1(\mathcal{L}(B)) \le d$ and $\lambda_1(\mathcal{L}(B)) > \gamma d$, under a promise. \\
			\hline
			
			\textbf{Lattice-based} 
			& CVP (Closest Vector Problem) 
			& Search 
			& Find the lattice point closest to a target 
			& Given a basis $B$ and target $\bm t$, output $\bm v \in \mathcal{L}(B)$ minimizing $\|\bm t-\bm v\|$. \\
			\hline
			
			\textbf{Lattice-based} 
			& Approx-CVP 
			& Search 
			& Find a lattice vector sufficiently close to the target 
			& Given $(B,\bm t)$, output $\bm v \in \mathcal{L}(B)$ with $\|\bm t-\bm v\| \le \gamma \cdot \mathrm{dist}(\bm t,\mathcal{L}(B))$. \\
			\hline
			
			\textbf{Lattice-based} 
			& GapCVP 
			& Decision 
			& Decide whether the target is very close or definitely far from the lattice 
			& Given $(B,\bm t,d)$, distinguish between $\mathrm{dist}(\bm t,\mathcal{L}(B)) \le d$ and $> \gamma d$, under a promise. \\
			\hline
			
			\textbf{Lattice-based} 
			& BDD (Bounded Distance Decoding) 
			& Search 
			& Recover the unique nearby lattice point 
			& Given $(B,\bm t)$ with promise that $\bm t$ is unusually close to some unique $\bm v \in \mathcal{L}(B)$, output that $\bm v$. \\
			\hline
			
			\textbf{Lattice-based} 
			& SIS (Short Integer Solution) 
			& Search 
			& Find a short nonzero modular relation 
			& Given $A \in \mathbb{Z}_q^{m\times n}$, find $\bm z \in \mathbb{Z}^n\setminus\{\bm 0\}$ such that $A\bm z \equiv \bm 0 \pmod q$ and $\|\bm z\|$ is small. \\
			\hline
			
			\textbf{Lattice-based} 
			& Decision-SIS style variant 
			& Decision 
			& Decide whether a short modular dependency exists 
			& Given $(A,\beta)$, decide whether there exists nonzero $\bm z$ with $A\bm z \equiv \bm 0 \pmod q$ and $\|\bm z\|\le \beta$. \\
			\hline
			
			\textbf{Lattice-based} 
			& Search-LWE 
			& Search 
			& Recover the hidden secret from noisy linear equations 
			& Given samples $(\bm a_i,b_i)$ where $b_i=\langle \bm a_i,\bm s\rangle+e_i \pmod q$, recover $\bm s$. \\
			\hline
			
			\textbf{Lattice-based} 
			& Decision-LWE 
			& Distinguishing 
			& Tell LWE samples from uniform random samples 
			& Distinguish $(\bm a,\langle \bm a,\bm s\rangle+e \bmod q)$ from $(\bm a,u)$ where $u \leftarrow \mathbb{Z}_q$ uniformly. \\
			\hline
			
			\textbf{Lattice-based} 
			& Search-RLWE / Search-MLWE 
			& Search 
			& Recover the secret ring/module element 
			& Given noisy ring or module linear samples, recover the hidden secret element/vector. \\
			\hline
			
			\textbf{Lattice-based} 
			& Decision-RLWE / Decision-MLWE 
			& Distinguishing 
			& Distinguish structured noisy algebraic samples from random ones 
			& Distinguish valid RLWE/MLWE samples from uniformly random samples over the same space. \\
			\hline
			
			\textbf{Code-based} 
			& Syndrome Decoding (SD) 
			& Search 
			& Find a low-weight error matching the syndrome 
			& Given $H,\bm s,w$, find $\bm e$ such that $H\bm e^T=\bm s^T$ and $w_H(\bm e)\le w$. \\
			\hline
			
			\textbf{Code-based} 
			& Decoding / Nearest Codeword Problem 
			& Search 
			& Find the codeword closest to a received word 
			& Given a code $C$ and received word $\bm y$, output $\bm c \in C$ minimizing Hamming distance $d_H(\bm y,\bm c)$. \\
			\hline
			
			\textbf{Code-based} 
			& Low-Weight Codeword Problem 
			& Search 
			& Find a nonzero codeword of small Hamming weight 
			& Given a linear code $C$, output nonzero $\bm c \in C$ with small $w_H(\bm c)$. \\
			\hline
			
			\textbf{Code-based} 
			& Decision-SD 
			& Decision 
			& Decide whether a low-weight error exists 
			& Given $(H,\bm s,w)$, decide whether there exists $\bm e$ such that $H\bm e^T=\bm s^T$ and $w_H(\bm e)\le w$. \\
			\hline
			
			\textbf{Code-based} 
			& Gap decoding variant 
			& Decision 
			& Decide whether decoding radius is small or definitely larger 
			& Given an instance and threshold $w$, distinguish between existence of a solution of weight $\le w$ and nonexistence of any solution of weight $\le \gamma w$. \\
			\hline
			
			\textbf{Code-based} 
			& Structured-code indistinguishability 
			& Distinguishing 
			& Tell whether a public code hides algebraic structure or is random 
			& Distinguish a disguised structured code family from a uniformly random linear code. \\
			\hline
			
			\textbf{Code-based} 
			& McEliece public-key indistinguishability viewpoint 
			& Distinguishing 
			& Distinguish a public generator/parity-check matrix from a random one 
			& Given a public matrix, decide whether it comes from a hidden efficiently decodable structured code or from a random code ensemble. \\
			\hline
			
			\textbf{Multivariate} 
			& MQ (Multivariate Quadratic solving) 
			& Search 
			& Find a solution to a system of multivariate quadratic equations 
			& Given quadratic polynomials $f_1,\dots,f_m$ over $\mathbb{F}_q$, find $\bm x$ such that $f_i(\bm x)=0$ for all $i$. \\
			\hline
			
			\textbf{Multivariate} 
			& Decision-MQ 
			& Decision 
			& Decide whether the polynomial system has a solution 
			& Given $f_1,\dots,f_m$, decide whether there exists $\bm x$ such that $f_i(\bm x)=0$ for every $i$. \\
			\hline
			
			\textbf{Multivariate} 
			& Oil/Vinegar structure distinguishability style question 
			& Distinguishing 
			& Tell whether a public polynomial map has hidden structure or is random-looking 
			& Distinguish a public key generated from a structured trapdoor construction from a uniformly random quadratic system. \\
			\hline
			
			\textbf{Isogeny-based} 
			& Path-finding / hidden isogeny recovery 
			& Search 
			& Recover the secret isogeny or path in the isogeny graph 
			& Given public elliptic-curve data related by an unknown isogeny, compute the connecting isogeny or secret path. \\
			\hline
			
			\textbf{Isogeny-based} 
			& Endomorphism / auxiliary hidden-structure recovery 
			& Search 
			& Recover hidden algebraic structure enabling the attack 
			& Given public objects, compute the hidden endomorphism/isogeny information needed to break the scheme. \\
			\hline
			
			\textbf{Isogeny-based} 
			& Decision supersingularity-type problems 
			& Decision 
			& Decide whether an elliptic curve has a target structural property 
			& Given a curve or related object, decide whether it belongs to a particular promised class, such as supersingular vs.\ ordinary. \\
			\hline
			
			\textbf{Isogeny-based} 
			& Isogeny-class / key-distribution distinguishability style question 
			& Distinguishing 
			& Tell whether public data comes from a structured isogeny process or random sampling 
			& Distinguish public data generated by the cryptosystem from randomly generated objects in the same ambient space. \\
			\hline
			
			\textbf{Hash-based} 
			& Second-preimage / preimage finding on hash trees or one-way functions 
			& Search 
			& Find a secret input consistent with a public digest 
			& Given output $y=f(x)$ or a hash digest, recover $x$ or another valid input mapping to the same target. \\
			\hline
			
			\textbf{Hash-based} 
			& Collision finding 
			& Search 
			& Find two distinct inputs with the same hash output 
			& Given a hash function $H$, output $x \neq x'$ such that $H(x)=H(x')$. \\
			\hline
			
			\textbf{Hash-based} 
			& One-wayness decision style variant 
			& Decision 
			& Decide whether a valid preimage under a bound or promise exists 
			& Given an output and promise constraints, decide whether there exists an input satisfying the specified relation. \\
			\hline
			
			\textbf{Hash-based} 
			& Pseudorandom output distinguishability style assumption 
			& Distinguishing 
			& Tell whether outputs behave like a random oracle / random function or a concrete structured primitive 
			& Distinguish outputs of a concrete construction from truly random outputs, when such a formulation is relevant. \\
			\hline
			
		\end{tabular}
	\end{center}
	\normalsize
	
	\begin{center}
		\small
		\renewcommand{\arraystretch}{1.35}
		\begin{tabular}{|p{2.2cm}|p{3.8cm}|p{2cm}|p{6.7cm}|}
			\hline
			\textbf{Family} & \textbf{Problem} & \textbf{Type} & \textbf{What the adversary must do} \\
			\hline
			
			Lattice-based & SVP, Approx-SVP, CVP, BDD, SIS, Search-LWE & Search & Recover a shortest/close vector, a short modular relation, or an LWE secret. \\
			\hline
			
			Lattice-based & GapSVP, GapCVP, decision-style SIS variants & Decision & Answer whether a short vector / close vector / short relation exists below a threshold. \\
			\hline
			
			Lattice-based & Decision-LWE, Decision-RLWE, Decision-MLWE & Distinguishing & Tell structured noisy samples from uniformly random samples. \\
			\hline
			
			Code-based & Syndrome Decoding, nearest codeword decoding, low-weight codeword search & Search & Recover a sparse error, decode a received word, or find a low-weight codeword. \\
			\hline
			
			Code-based & Decision-SD, gap decoding variants & Decision & Decide whether a low-weight error or nearby codeword exists. \\
			\hline
			
			Code-based & Structured code vs.\ random code, McEliece public-key indistinguishability & Distinguishing & Tell hidden algebraic code structure from a random linear code. \\
			\hline
			
			Multivariate & MQ solving & Search & Find a solution to a system of multivariate quadratic equations. \\
			\hline
			
			Multivariate & Decision-MQ & Decision & Decide whether the quadratic system has any solution. \\
			\hline
			
			Multivariate & Structured public map vs.\ random quadratic map & Distinguishing & Tell whether the public polynomial system hides trapdoor structure. \\
			\hline
			
			Isogeny-based & Secret isogeny / path recovery & Search & Recover the hidden isogeny or secret path in the graph. \\
			\hline
			
			Isogeny-based & Supersingularity-type membership questions & Decision & Decide whether a given object has a target structural property. \\
			\hline
			
			Isogeny-based & Structured public data vs.\ random object distributions & Distinguishing & Tell cryptosystem-generated objects from random ambient objects. \\
			\hline
			
			Hash-based & Preimage, second-preimage, collision finding & Search & Find an input or pair of inputs satisfying the required hash relation. \\
			\hline
			
			Hash-based & Promise-style existence questions & Decision & Decide whether a valid witness exists under the given promise. \\
			\hline
			
			Hash-based & Random-oracle / pseudorandom-style distinguishability views & Distinguishing & Tell whether outputs come from a specific construction or an ideal random source. \\
			\hline
			
		\end{tabular}
	\end{center}
	
	\paragraph{Reading guide.}
	A \emph{search} problem asks the adversary to recover an explicit hidden object.
	A \emph{decision} problem asks only for a yes/no answer about whether an instance has a target property.
	A \emph{distinguishing} problem asks the adversary to tell apart two distributions, usually a structured distribution and a random one.
	
	\newpage
	\section{PQC hard problems classified as Search, Decision, and Distinguishing}
	
	\begin{center}
		\scriptsize
		\renewcommand{\arraystretch}{1.35}
		\setlength{\tabcolsep}{4pt}
		\begin{tabular}{|p{2.1cm}|p{3.2cm}|p{1.8cm}|p{3.8cm}|p{5.0cm}|}
			\hline
			\textbf{PQC family} & \textbf{Representative hard problem / assumption} & \textbf{Type} & \textbf{Informal task} & \textbf{Typical formulation} \\
			\hline
			
			\multirow{12}{=}{\textbf{Lattice-based}}
			& SVP (Shortest Vector Problem)
			& Search
			& Find the shortest nonzero lattice vector
			& Given a basis $B$ of a lattice $\mathcal{L}(B)$, output $\bm v \in \mathcal{L}(B)\setminus\{\bm 0\}$ such that $\|\bm v\|=\lambda_1(\mathcal{L}(B))$. \\
			\cline{2-5}
			
			& Approx-SVP
			& Search
			& Find a sufficiently short nonzero lattice vector
			& Given $B$, output $\bm v \neq \bm 0$ with $\|\bm v\| \le \gamma \lambda_1(\mathcal{L}(B))$. \\
			\cline{2-5}
			
			& GapSVP
			& Decision
			& Decide whether the lattice contains a very short vector
			& Given $(B,d)$, distinguish between $\lambda_1(\mathcal{L}(B)) \le d$ and $\lambda_1(\mathcal{L}(B)) > \gamma d$, under a promise. \\
			\cline{2-5}
			
			& CVP (Closest Vector Problem)
			& Search
			& Find the lattice vector closest to a target
			& Given a basis $B$ and target $\bm t$, output $\bm v \in \mathcal{L}(B)$ minimizing $\|\bm t-\bm v\|$. \\
			\cline{2-5}
			
			& Approx-CVP
			& Search
			& Find a lattice vector sufficiently close to the target
			& Given $(B,\bm t)$, output $\bm v \in \mathcal{L}(B)$ with $\|\bm t-\bm v\| \le \gamma \cdot \mathrm{dist}(\bm t,\mathcal{L}(B))$. \\
			\cline{2-5}
			
			& GapCVP
			& Decision
			& Decide whether the target is very close or definitely far
			& Given $(B,\bm t,d)$, distinguish between $\mathrm{dist}(\bm t,\mathcal{L}(B)) \le d$ and $> \gamma d$, under a promise. \\
			\cline{2-5}
			
			& BDD (Bounded Distance Decoding)
			& Search
			& Recover the unique nearby lattice point
			& Given $(B,\bm t)$ with a uniqueness promise, output the lattice vector closest to $\bm t$. \\
			\cline{2-5}
			
			& SIS (Short Integer Solution)
			& Search
			& Find a short nonzero modular relation
			& Given $A \in \mathbb{Z}_q^{m\times n}$, find $\bm z \in \mathbb{Z}^n\setminus\{\bm 0\}$ such that $A\bm z \equiv \bm 0 \pmod q$ and $\|\bm z\|$ is small. \\
			\cline{2-5}
			
			& Decision-SIS style variant
			& Decision
			& Decide whether a short modular dependency exists
			& Given $(A,\beta)$, decide whether there exists nonzero $\bm z$ with $A\bm z \equiv \bm 0 \pmod q$ and $\|\bm z\|\le \beta$. \\
			\cline{2-5}
			
			& Search-LWE
			& Search
			& Recover the hidden secret from noisy linear equations
			& Given samples $(\bm a_i,b_i)$ where $b_i=\langle \bm a_i,\bm s\rangle+e_i \pmod q$, recover $\bm s$. \\
			\cline{2-5}
			
			& Decision-LWE
			& Distinguishing
			& Tell LWE samples from uniform random samples
			& Distinguish $(\bm a,\langle \bm a,\bm s\rangle+e \bmod q)$ from $(\bm a,u)$ where $u \leftarrow \mathbb{Z}_q$ uniformly. \\
			\cline{2-5}
			
			& Decision-RLWE / Decision-MLWE
			& Distinguishing
			& Distinguish structured noisy algebraic samples from random ones
			& Distinguish valid RLWE/MLWE samples from uniformly random samples over the same space. \\
			\hline
			
			\multirow{6}{=}{\textbf{Code-based}}
			& Syndrome Decoding (SD)
			& Search
			& Find a low-weight error matching the syndrome
			& Given $H,\bm s,w$, find $\bm e$ such that $H\bm e^T=\bm s^T$ and $w_H(\bm e)\le w$. \\
			\cline{2-5}
			
%			& Nearest Codeword Problem
%			& Search
%			& Find the codeword closest to a received word
%			& Given a code $C$ and received word $\bm y$, output $\bm c \in C$ minimizing Hamming distance $d_H(\bm y,\bm c)$. \\
%			\cline{2-5}
%			
%			& Low-Weight Codeword Problem
%			& Search
%			& Find a nonzero codeword of small Hamming weight
%			& Given a linear code $C$, output nonzero $\bm c \in C$ with small $w_H(\bm c)$. \\
%			\cline{2-5}
%			
%			& Decision-SD
%			& Decision
%			& Decide whether a low-weight error exists
%			& Given $(H,\bm s,w)$, decide whether there exists $\bm e$ such that $H\bm e^T=\bm s^T$ and $w_H(\bm e)\le w$. \\
%			\cline{2-5}
%			
%			& Gap decoding variant
%			& Decision
%			& Decide whether decoding radius is small or definitely larger
%			& Distinguish between existence of a solution of weight $\le w$ and nonexistence of any solution of weight $\le \gamma w$. \\
%			\cline{2-5}
%			
%			& Structured-code indistinguishability
%			& Distinguishing
%			& Tell whether a public code hides algebraic structure or is random
%			& Distinguish a disguised structured code family from a uniformly random linear code. \\
%			\hline
%			
%			\multirow{3}{=}{\textbf{Multivariate}}
%			& MQ (Multivariate Quadratic solving)
%			& Search
%			& Find a solution to a quadratic polynomial system
%			& Given quadratic polynomials $f_1,\dots,f_m$ over $\mathbb{F}_q$, find $\bm x$ such that $f_i(\bm x)=0$ for all $i$. \\
%			\cline{2-5}
%			
%			& Decision-MQ
%			& Decision
%			& Decide whether the polynomial system has a solution
%			& Given $f_1,\dots,f_m$, decide whether there exists $\bm x$ such that $f_i(\bm x)=0$ for every $i$. \\
%			\cline{2-5}
%			
%			& Structured public map indistinguishability
%			& Distinguishing
%			& Tell whether the public key hides trapdoor structure
%			& Distinguish a structured public quadratic map from a uniformly random-looking quadratic system. \\
%			\hline
%			
%			\multirow{3}{=}{\textbf{Isogeny-based}}
%			& Hidden isogeny / path recovery
%			& Search
%			& Recover the secret isogeny or path
%			& Given public elliptic-curve data related by an unknown isogeny, compute the connecting isogeny or secret path. \\
%			\cline{2-5}
%			
%			& Structural membership problem
%			& Decision
%			& Decide whether an object has a target property
%			& Given a curve or related object, decide whether it belongs to a promised class, such as supersingular vs.\ ordinary. \\
%			\cline{2-5}
%			
%			& Public-data distinguishability
%			& Distinguishing
%			& Tell whether public data comes from a cryptosystem or a random source
%			& Distinguish public data generated by a structured isogeny process from random ambient objects. \\
%			\hline
%			
%			\multirow{3}{=}{\textbf{Hash-based}}
%			& Preimage / second-preimage / collision finding
%			& Search
%			& Find an input or pair of inputs satisfying a hash relation
%			& Given a hash function $H$, recover a preimage, second preimage, or collision depending on the task. \\
%			\cline{2-5}
%			
%			& Promise-style existence variant
%			& Decision
%			& Decide whether a valid witness exists under a promise
%			& Given an output and constraints, decide whether there exists an input satisfying the specified relation. \\
%			\cline{2-5}
%			
%			& Pseudorandom-output distinguishability style view
%			& Distinguishing
%			& Tell whether outputs come from a real construction or an ideal random source
%			& Distinguish outputs of a concrete construction from truly random outputs when modeled that way. \\
%			\hline
		\end{tabular}
	\end{center}
	\normalsize
	
	\newpage
	\section{PQC hard problems classified as Search, Decision, and Distinguishing}
	
	\begin{center}
		\scriptsize
		\renewcommand{\arraystretch}{1.35}
		\setlength{\tabcolsep}{4pt}
		\begin{tabular}{|p{2.1cm}|p{1.8cm}|p{3.3cm}|p{3.8cm}|p{4.9cm}|}
			\hline
			\textbf{PQC family} & \textbf{Type} & \textbf{Representative hard problem / assumption} & \textbf{Informal task} & \textbf{Typical formulation} \\
			\hline
			
			\multirow{12}{=}{\textbf{Lattice-based}}
			& \multirow{7}{=}{Search}
			& SVP (Shortest Vector Problem)
			& Find the shortest nonzero lattice vector
			& Given a basis $B$ of a lattice $\mathcal{L}(B)$, output $\bm v \in \mathcal{L}(B)\setminus\{\bm 0\}$ such that $\|\bm v\|=\lambda_1(\mathcal{L}(B))$. \\
			\cline{3-5}
			
			& & Approx-SVP
			& Find a sufficiently short nonzero lattice vector
			& Given $B$, output $\bm v \neq \bm 0$ with $\|\bm v\| \le \gamma \lambda_1(\mathcal{L}(B))$. \\
			\cline{3-5}
			
			& & CVP (Closest Vector Problem)
			& Find the lattice vector closest to a target
			& Given a basis $B$ and target $\bm t$, output $\bm v \in \mathcal{L}(B)$ minimizing $\|\bm t-\bm v\|$. \\
			\cline{3-5}
			
			& & Approx-CVP
			& Find a lattice vector sufficiently close to the target
			& Given $(B,\bm t)$, output $\bm v \in \mathcal{L}(B)$ with $\|\bm t-\bm v\| \le \gamma \cdot \mathrm{dist}(\bm t,\mathcal{L}(B))$. \\
			\cline{3-5}
			
			& & BDD (Bounded Distance Decoding)
			& Recover the unique nearby lattice point
			& Given $(B,\bm t)$ with a uniqueness promise, output the lattice vector closest to $\bm t$. \\
			\cline{3-5}
			
			& & SIS (Short Integer Solution)
			& Find a short nonzero modular relation
			& Given $A \in \mathbb{Z}_q^{m\times n}$, find $\bm z \in \mathbb{Z}^n\setminus\{\bm 0\}$ such that $A\bm z \equiv \bm 0 \pmod q$ and $\|\bm z\|$ is small. \\
			\cline{3-5}
			
			& & Search-LWE
			& Recover the hidden secret from noisy linear equations
			& Given samples $(\bm a_i,b_i)$ where $b_i=\langle \bm a_i,\bm s\rangle+e_i \pmod q$, recover $\bm s$. \\
			\cline{2-5}
			
			& \multirow{3}{=}{Decision}
			& GapSVP
			& Decide whether the lattice contains a very short vector
			& Given $(B,d)$, distinguish between $\lambda_1(\mathcal{L}(B)) \le d$ and $\lambda_1(\mathcal{L}(B)) > \gamma d$, under a promise. \\
			\cline{3-5}
			
			& & GapCVP
			& Decide whether the target is very close or definitely far
			& Given $(B,\bm t,d)$, distinguish between $\mathrm{dist}(\bm t,\mathcal{L}(B)) \le d$ and $> \gamma d$, under a promise. \\
			\cline{3-5}
			
			& & Decision-SIS style variant
			& Decide whether a short modular dependency exists
			& Given $(A,\beta)$, decide whether there exists nonzero $\bm z$ with $A\bm z \equiv \bm 0 \pmod q$ and $\|\bm z\|\le \beta$. \\
			\cline{2-5}
			
			& \multirow{2}{=}{Distinguishing}
			& Decision-LWE
			& Tell LWE samples from uniform random samples
			& Distinguish $(\bm a,\langle \bm a,\bm s\rangle+e \bmod q)$ from $(\bm a,u)$ where $u \leftarrow \mathbb{Z}_q$ uniformly. \\
			\cline{3-5}
			
			& & Decision-RLWE / Decision-MLWE
			& Distinguish structured noisy algebraic samples from random ones
			& Distinguish valid RLWE/MLWE samples from uniformly random samples over the same space. \\
			\hline
			
			\multirow{6}{=}{\textbf{Code-based}}
			& \multirow{3}{=}{Search}
			& Syndrome Decoding (SD)
			& Find a low-weight error matching the syndrome
			& Given $H,\bm s,w$, find $\bm e$ such that $H\bm e^T=\bm s^T$ and $w_H(\bm e)\le w$. \\
			\cline{3-5}
			
			& & Nearest Codeword Problem
			& Find the codeword closest to a received word
			& Given a code $C$ and received word $\bm y$, output $\bm c \in C$ minimizing Hamming distance $d_H(\bm y,\bm c)$. \\
			\cline{3-5}
			
			& & Low-Weight Codeword Problem
			& Find a nonzero codeword of small Hamming weight
			& Given a linear code $C$, output nonzero $\bm c \in C$ with small $w_H(\bm c)$. \\
			\cline{2-5}
			
			& \multirow{2}{=}{Decision}
			& Decision-SD
			& Decide whether a low-weight error exists
			& Given $(H,\bm s,w)$, decide whether there exists $\bm e$ such that $H\bm e^T=\bm s^T$ and $w_H(\bm e)\le w$. \\
			\cline{3-5}
			
			& & Gap decoding variant
			& Decide whether decoding radius is small or definitely larger
			& Distinguish between existence of a solution of weight $\le w$ and nonexistence of any solution of weight $\le \gamma w$. \\
			\cline{2-5}
			
			& Distinguishing
			& Structured-code indistinguishability
			& Tell whether a public code hides algebraic structure or is random
			& Distinguish a disguised structured code family from a uniformly random linear code. \\
			\hline
			
%			\multirow{3}{=}{\textbf{Multivariate}}
%			& Search
%			& MQ (Multivariate Quadratic solving)
%			& Find a solution to a quadratic polynomial system
%			& Given quadratic polynomials $f_1,\dots,f_m$ over $\mathbb{F}_q$, find $\bm x$ such that $f_i(\bm x)=0$ for all $i$. \\
%			\cline{2-5}
%			
%			& Decision
%			& Decision-MQ
%			& Decide whether the polynomial system has a solution
%			& Given $f_1,\dots,f_m$, decide whether there exists $\bm x$ such that $f_i(\bm x)=0$ for every $i$. \\
%			\cline{2-5}
%			
%			& Distinguishing
%			& Structured public map indistinguishability
%			& Tell whether the public key hides trapdoor structure
%			& Distinguish a structured public quadratic map from a uniformly random-looking quadratic system. \\
%			\hline
%			
%			\multirow{3}{=}{\textbf{Isogeny-based}}
%			& Search
%			& Hidden isogeny / path recovery
%			& Recover the secret isogeny or path
%			& Given public elliptic-curve data related by an unknown isogeny, compute the connecting isogeny or secret path. \\
%			\cline{2-5}
%			
%			& Decision
%			& Structural membership problem
%			& Decide whether an object has a target property
%			& Given a curve or related object, decide whether it belongs to a promised class, such as supersingular vs.\ ordinary. \\
%			\cline{2-5}
%			
%			& Distinguishing
%			& Public-data distinguishability
%			& Tell whether public data comes from a cryptosystem or a random source
%			& Distinguish public data generated by a structured isogeny process from random ambient objects. \\
%			\hline
			
%			\multirow{3}{=}{\textbf{Hash-based}}
%			& Search
%			& Preimage / second-preimage / collision finding
%			& Find an input or pair of inputs satisfying a hash relation
%			& Given a hash function $H$, recover a preimage, second preimage, or collision depending on the task. \\
%			\cline{2-5}
%			
%			& Decision
%			& Promise-style existence variant
%			& Decide whether a valid witness exists under a promise
%			& Given an output and constraints, decide whether there exists an input satisfying the specified relation. \\
%			\cline{2-5}
%			
%			& Distinguishing
%			& Pseudorandom-output distinguishability style view
%			& Tell whether outputs come from a real construction or an ideal random source
%			& Distinguish outputs of a concrete construction from truly random outputs when modeled that way. \\
%			\hline
			
		\end{tabular}
	\end{center}
	\normalsize
	
\end{document}